% Generated draft: requires mathematical and source review.
\documentclass[11pt]{article}
\usepackage[margin=25mm]{geometry}
\usepackage{fontspec}
\setmainfont{Latin Modern Roman}
\usepackage{amsmath,amssymb,mathtools,unicode-math}
\setmathfont{Latin Modern Math}
\usepackage{enumitem,needspace,xurl,hyperref}
\setlength{\parindent}{0pt}
\setlength{\parskip}{5pt}
\setlength{\emergencystretch}{3em}
\newcommand{\N}{\mathbb{N}}
\newcommand{\Z}{\mathbb{Z}}
\newcommand{\Q}{\mathbb{Q}}
\newcommand{\R}{\mathbb{R}}
\newcommand{\C}{\mathbb{C}}
\newcommand{\F}{\mathbb{F}}
\newcommand{\A}{\mathbb{A}}
\newcommand{\PP}{\mathbb{P}}
\newcommand{\E}{\mathbb{E}}
\newcommand{\eps}{\varepsilon}
\newcommand{\dd}{\,\mathrm{d}}
\begin{document}
\section*{UnsolvedMath ID: 701099}
\textbf{Source problem code:} catalogue entry
\subsection{Definitions and mathematical statement}
For a group \(G\), define its group of power automorphisms by
\[
\operatorname{PAut}(G)=\{\alpha\in\operatorname{Aut}(G):\alpha(H)=H\text{ for every subgroup }H\leq G\},
\]
with composition as the group operation. Thus invariance means setwise invariance, not pointwise fixation.

Is it true that for every prime \(p\) and every finite abelian group \(A\) of order \(p^n\), where \(n\geq0\) is an integer, there exists a group \(G\) such that
\[
\operatorname{PAut}(G)\cong A?
\]
No finiteness restriction on \(G\) is stated in the supplied source. The case \(n=0\) is included under the usual convention that the trivial group is a finite abelian \(p\)-group. The statement bears the attribution ``M. Morigi'', while the supplied proposer metadata names V. G. Bardakov; their respective roles are not explained.
\subsection{Short English statement}
Does every finite abelian \(p\)-group occur, up to isomorphism, as the full group of automorphisms preserving every subgroup of some group? The realizing group is not required to be finite in the supplied statement.
\subsection{Sources}
\begin{itemize}
\item Local UnsolvedMath registry, permanent ID \(701099\), exact source code \texttt{catalogue entry}, title \emph{Kourovka Notebook Problem 15.11}; supplied upstream revision \texttt{d371ee917ac9647fbb06cd6fc529deec6f3dda82}.
\item Supplied source attribution: \emph{The Kourovka Notebook (Unsolved Problems in Group Theory)}, original problem 15.11, 15th issue (2002), as recorded with the 21st issue edition (2026), arXiv:1401.0300v48. The collection metadata names E. I. Khukhro and V. D. Mazurov as editors.
\item The statement is prefixed by ``M. Morigi''; the background and proposer metadata name V. G. Bardakov, with proposed year 2002. The supplied record does not explain this attribution discrepancy.
\item The local record labels the problem \emph{open} and records an update date of 2026-10-16. These are supplied registry metadata, not an independently verified assessment of the literature or mathematical status. The cited source text has not been independently consulted for this draft.
\end{itemize}
\end{document}
